\section{Rings, fields, and the prime field}
\label{sec:finite-field-foundations}

Changing a modulus changes the rules of division. In $\Z/7\Z$ every
nonzero residue has an inverse; in $\Z/8\Z$, $2\cdot4=0$ although
neither factor is zero. We begin by naming the structures that distinguish
these examples.

\begin{definition}
A \emph{ring} has an abelian group operation called addition, an associative
multiplication with identity $1\ne0$, and both distributive laws
\[
a(b+c)=ab+ac,\qquad (a+b)c=ac+bc.
\]
Multiplication need not commute. A \emph{unit} is an element with a
two-sided multiplicative inverse. A \emph{division ring} is a ring in
which every nonzero element is a unit. A \emph{field} is a commutative
division ring. A nonzero element $a$ is a \emph{zero divisor} if
$ab=0$ or $ba=0$ for some nonzero $b$.
\end{definition}

The integers are a commutative ring, but not a field. The rational numbers
are a field. Residues modulo $n$ form a commutative ring whose units are
exactly the classes coprime to $n$. Matrices give familiar noncommutative
rings: matrix multiplication distributes over addition but generally
$AB\ne BA$. A unit cannot be a zero divisor, because multiplying by its
inverse would force the other factor to be zero.

\begin{proposition}
A finite commutative ring with no zero divisors is a field. In particular,
for $n\ge2$, $\Z/n\Z$ is a field exactly when $n$ is prime.
\end{proposition}
\begin{proof}
Fix a nonzero $a$. If $ax=ay$, then $a(x-y)=0$, so $x=y$.
Multiplication by $a$ is therefore injective on a finite set and hence
surjective. In particular $ab=1$ for some $b$, which is an inverse by
commutativity. For prime $p$, Euclid's lemma excludes zero divisors modulo
$p$; alternatively, B\'ezout gives each inverse. For composite $n=ab$,
with $1<a,b<n$, the nonzero residues of $a,b$ have product zero.
\end{proof}

\begin{definition}
The \emph{characteristic} of a ring is the least positive integer $m$
such that $m\cdot1=0$, if such an integer exists; otherwise it is $0$.
Here $m\cdot1$ means repeated addition, not exponentiation.
\end{definition}

\begin{proposition}
Every finite field has prime characteristic $p$ and contains a distinguished
copy of $\F_p$, called its \emph{prime field}.
\end{proposition}
\begin{proof}
Finiteness makes two of $0,1,1+1,\ldots$ equal, so a positive multiple
of $1$ is zero. Let $m$ be the least such multiple. If $m=ab$ with
$1<a,b<m$, then $(a\cdot1)(b\cdot1)=0$ has two nonzero factors,
impossible in a field. Also $m\ne1$, so $m=p$ is prime. The map
$\Z\to K$, $a\mapsto a\cdot1$, respects both operations and has kernel
$p\Z$. Thus its image is a subfield identified with $\Z/p\Z=\F_p$.
\end{proof}

\section{Coordinates, exchange, and dimension}
\label{sec:finite-field-dimension}

A vector space over a field $k$ is an abelian group $(V,+)$ with scalar
multiplication satisfying
\[
a(u+v)=au+av,\quad (a+b)v=av+bv,\quad
(ab)v=a(bv),\quad 1v=v.
\]
The scalars belong to $k$, and the vectors belong to $V$. For example,
in $\F_3^2$, $2(1,2)=(2,1)$. If a field $K$ contains $k$, its own
elements form a vector space over $k$: scalar multiplication is the
field multiplication restricted to scalars in $k$.

\begin{definition}
The \emph{span} of $v_1,\ldots,v_r$ is the set of their linear
combinations $\sum a_iv_i$. The list is \emph{linearly independent} if
$\sum a_iv_i=0$ implies every $a_i=0$. A \emph{basis} is an independent
spanning list. Its coordinates express every vector uniquely:
spanning gives existence, and subtracting two expressions and using
independence gives uniqueness.
\end{definition}

An empty list spans $\{0\}$. If an independent list does not yet span,
choose $w$ outside its span. The enlarged list remains independent:
in a relation $\sum a_iv_i+bw=0$, a nonzero $b$ would put $w$ in the old
span. Thus $b=0$, and all the $a_i$ also vanish. In a finite vector
space this construction stops, because each span strictly enlarges.
It produces a basis.

Before defining dimension we must know that different bases have the
same length. Counting coordinates will prove this for finite fields,
but the following argument works over any field.

\begin{lemma}[Exchange]
If $v_1,\ldots,v_r$ are independent and $w_1,\ldots,w_m$ span $V$,
then $r\le m$.
\end{lemma}
\begin{proof}
Insert the $v_i$ one at a time into the spanning list. After $k$
insertions, suppose the list consists of $v_1,\ldots,v_k$ and $m-k$
old vectors and still spans. Express $v_{k+1}$ in this list. Some old
vector has nonzero coefficient; otherwise $v_{k+1}$ is in the span of
the previous $v_i$, contradicting independence. Solve for that old
vector by dividing by its coefficient. Replacing it by $v_{k+1}$
preserves spanning, since the solved equation recovers the discarded
vector. Previously inserted vectors are retained.

If $r>m$, after $m$ insertions $v_1,\ldots,v_m$ would span, forcing
$v_{m+1}$ into their span, again contradicting independence. Thus
$r\le m$. If $m=0$, $V=\{0\}$ has no nonempty independent list,
so the conclusion also holds.
\end{proof}

\begin{corollary}
Any two finite bases of a vector space have the same length. This length
is its \emph{dimension} over the scalar field, denoted $\dim_kV$.
\end{corollary}
\begin{proof}
Use each basis as the independent list and the other as the spanning
list. The two inequalities force equality.
\end{proof}

\begin{example}[Coordinates over a prime field]
In $\F_3^2$, put $u=(1,1)$ and $v=(1,2)$. The equations
$au+bv=0$ give $a+b=0$ and $a+2b=0$, hence $a=b=0$.
For an arbitrary $(x,y)$, the same equations with right sides $x,y$
give $b=y-x$ and $a=2x-y$. Thus $u,v$ are a basis, and
$(2,0)=u+v$ has coordinates $(1,1)$ in this basis.
The redundant list $u,v,u+v$ still spans but is not independent.
\end{example}

\begin{theorem}[The necessary size of a finite field]
\label{thm:finite-field-size}
A $d$-dimensional vector space over a field with $q$ elements has $q^d$
elements. In particular, every finite field has $p^d$ elements for a
prime $p$ and integer $d\ge1$.
\end{theorem}
\begin{proof}
Choose a basis. Each of its $d$ coordinates has $q$ choices, and the
coordinate map $k^d\to V$ is bijective by existence and uniqueness.
For a finite field $K$, use its prime field $\F_p$ as scalars. The
basis-building procedure above stops because $K$ is finite, so
$|K|=p^d$. Since $1\ne0$, $d\ge1$.
\end{proof}

Thus there is no field with $6$ or $12$ elements. A field with $81$
elements must have characteristic $3$ and dimension $4$ over $\F_3$.
At this point prime-power size is only a necessary condition. We will
construct every prime-power size and prove uniqueness in
Section~\ref{sec:finite-field-classification}. First, a quadratic example
makes the arithmetic concrete.
